% !PRINTED 757
\weberpage{757}

Let us therefore assume that $\pi$, and consequently also $i\pi$, is an algebraic number; then $i\pi$ is one of the roots of an irreducible rational equation $\chi(x)=0$ whose coefficients are integral rational numbers.

If we denote all the roots of this algebraic equation by $\beta_1,\beta_2,\ldots,\beta_\nu$, and denote the coefficient of $x^\nu$ in $\chi$ by $a$, then
\weq{2}{\chi(x) = a(x-\beta_1)(x-\beta_2)\ldots(x-\beta_\nu),}
and the products $a\beta_1,a\beta_2,\ldots,a\beta_\nu$ are \gesperrt{integral} algebraic numbers, and hence their integral symmetric functions are \gesperrt{integral rational numbers} (\S.\,133).

Since the number $i\pi$ is to occur among the $\beta$, we have by (1):
\[
(1+e^{\beta_1})(1+e^{\beta_2})\ldots(1+e^{\beta_\nu}) = 0,
\]
and if we carry out the multiplication, an equation results:
\[
1 + \sum e^{\beta_i} + \sum e^{\beta_i+\beta_h} + \sum e^{\beta_i+\beta_h+\beta_k} + \cdots = 0.
\]
Among the exponents in this sum the number zero may occur several times; let us assume $(C-1)$ times, so that $C$ is a positive integer, at least $=1$. The remaining exponents $\beta_i,\beta_i+\beta_h,\beta_i+\beta_h+\beta_k,\ldots$, some of which may also be equal to one another, we shall denote by $\alpha_1,\alpha_2,\ldots,\alpha_\mu$, so that the equation holds:
\weq{3}{C + e^{\alpha_1} + e^{\alpha_2} + \cdots + e^{\alpha_\mu} = 0.}
Here the $\alpha_1,\alpha_2,\ldots,\alpha_\mu$ are algebraic numbers which, multiplied by the integral rational number $a$, pass over into integral algebraic numbers.

But the symmetric functions of all the sums $\beta_i,\beta_i+\beta_h,\beta_i+\beta_h+\beta_k,\ldots$ are at the same time symmetric functions of the $\beta_1,\ldots,\beta_\nu$, and consequently rational numbers. These sums are therefore also the roots of a rational equation, and since the root $0$ can be separated off as often as it is present, the $\alpha_1,\alpha_2,\ldots,\alpha_\mu$ are likewise the roots of a rational equation; the symmetric fundamental functions of the $a\alpha_1,a\alpha_2,\ldots,a\alpha_\mu$ are \gesperrt{integral rational numbers}.

The absolute values of the numbers
\[
\alpha_1,\alpha_2,\ldots,\alpha_\mu
\]
shall now be denoted by
\[
a_1,a_2,\ldots,a_\mu
\]
respectively.
